\chapter{Fourier analysis and incidence of complex lines}
\label{chap:p06}

Let $D\subset\mathbb C^n$ be open, with $n\geq2$. A function on $D$ that is subharmonic on every complex-line section is bounded above on compact subsets of each individual line. The bounds could still depend on the line. To control them uniformly, we will place a probability measure on an annular set of high values, bound its energy independently of the function, and project it onto complex hyperplanes. Fourier analysis turns that energy bound into an $L^2$ bound for the projected densities. This forces one affine complex line to meet infinitely many high-value sets and contradicts its local upper bound.

The same Fourier transform also controls weak second derivatives. Its Hessian multipliers are bounded on every $L^p$, $1<p<\infty$, by a dyadic decomposition and cancellation. This supplies the real regularity needed to interpret determinants for nonsmooth obstacles later in the book.

\section{Fourier transform and energy}
\label{sec:fourier-energy}
Let $d\geq1$. A finite complex measure means $\mu=\sigma_1+\mathrm i\sigma_2$, with $\sigma_1,\sigma_2$ finite real signed measures. For $f\in L^1(\mathbb R^d)$ and such a measure $\mu$, use the Fourier transform\index{Fourier transform} convention
\[
 \widehat f(\xi)=\int_{\mathbb R^d}f(x)\mathrm e^{-2\pi\mathrm i x\cdot\xi}\,\mathrm d x,
 \qquad \widehat\mu(\xi)=\int\mathrm e^{-2\pi\mathrm i x\cdot\xi}\,\mathrm d\mu(x).
\]
These functions are bounded and continuous by dominated convergence. The Schwartz space\index{Schwartz space} $\mathcal S(\mathbb R^d)$ consists of smooth complex functions for which $x^\alpha\partial^\beta f$ is bounded for all multi-indices $\alpha,\beta$. In particular every derivative is integrable, and $C_c^\infty\subset\mathcal S$.

\begin{lemma}\label{lem:fourier-gaussian}
The Fourier transform maps $\mathcal S$ into itself. For $t>0$,
\[
 \widehat{\mathrm e^{-\pi t|x|^2}}(\xi)=t^{-d/2}\mathrm e^{-\pi|\xi|^2/t}.
\]
\end{lemma}
\begin{proof}
Differentiating under the integral multiplies $f$ by $-2\pi\mathrm i x_j$, and integration by parts multiplies its transform by $2\pi\mathrm i\xi_j$. Boundary terms vanish by rapid decay, first on boxes and then on letting their sides tend to infinity. Iterating gives boundedness of every $\xi^\alpha\partial^\beta\widehat f$, since it is a constant multiple of the transform of an integrable derivative of $x^\beta f$. This proves the Schwartz assertion.

In one dimension set $F(\xi)=\int\mathrm e^{-\pi x^2}\mathrm e^{-2\pi\mathrm i x\xi}\,\mathrm d x$. Integration by parts gives $F'=-2\pi\xi F$. Also $F(0)=1$: square the positive Gaussian integral and use planar polar coordinates to obtain $2\pi\int_0^\infty r\mathrm e^{-\pi r^2}\,\mathrm d r=1$. The differential equation gives $F(\xi)=\mathrm e^{-\pi\xi^2}$. Scaling and the product formula from Fubini give the result in dimension $d$.
\end{proof}

\begin{theorem}[Plancherel]\label{thm:plancherel}
\index{Plancherel's theorem}For $f\in\mathcal S$ the inversion formula is
\[
 f(x)=\int\widehat f(\xi)\mathrm e^{2\pi\mathrm i x\cdot\xi}\,\mathrm d\xi.
\]
The Fourier transform extends uniquely to a unitary operator on $L^2(\mathbb R^d)$, with inverse given on Schwartz functions by changing the exponential's sign. It agrees almost everywhere with the integral transform on $L^1\cap L^2$.
\end{theorem}
\begin{proof}
Insert the factor $\mathrm e^{-\pi\varepsilon^2|\xi|^2}$ into the frequency integrand in the asserted inversion formula. Fubini and \cref{lem:fourier-gaussian} identify it with $f*g_\varepsilon(x)$, where $g_\varepsilon(y)=\varepsilon^{-d}\mathrm e^{-\pi|y|^2/\varepsilon^2}$. This probability kernel is an approximate identity: uniform continuity controls the integral on $|y|<\delta$, and its mass on $|y|\geq\delta$ tends to zero by scaling and integrability. Thus $f*g_\varepsilon\to f$ uniformly. On the Fourier side, $\widehat f\in L^1$ permits dominated convergence, proving inversion. The same proof with the signs reversed gives inversion in the other order.

For $f,g\in\mathcal S$, apply inversion to $g$ and Fubini to get
\[
 \int f\overline g\,\mathrm d x=\int\widehat f\,\overline{\widehat g}\,\mathrm d\xi.
\]
The absolute integral needed for this interchange is at most $\|f\|_1\|\widehat g\|_1$. Smooth compactly supported functions are dense in $L^2$ by \cref{lem:mollifier-density}. Hence the identity extends the transform as an isometry to $L^2$. The inverse transform extends as an isometry too, and the two compositions are the identity on a dense subspace and therefore everywhere. This proves unitarity and uniqueness.

If $f\in L^1\cap L^2$, cutoff and smooth approximation give $f_j\in C_c^\infty$ converging in both norms. This follows by cutting off the tails, truncating the values, and using a mollifier; translation continuity in both norms gives the final approximation. The integral transforms converge uniformly by the $L^1$ estimate, while their $L^2$ transforms converge in $L^2$. A subsequence of an $L^2$ convergent sequence converges almost everywhere: select squared errors with summable integrals and apply Tonelli. The two limits agree almost everywhere.
\end{proof}

\begin{proposition}\label{prop:measure-fourier-l2}
If a finite positive measure $\mu$ on $\mathbb R^d$ satisfies $\widehat\mu\in L^2$, then $\mu=g\,\mathrm d x$ for a nonnegative $g\in L^1\cap L^2$, and $\|g\|_2=\|\widehat\mu\|_2$.
\end{proposition}
\begin{proof}
Let $g$ be the inverse $L^2$ transform of $\widehat\mu$. For $\phi\in\mathcal S$, inversion and Fubini give
\[
 \int\phi\,\mathrm d\mu=\int\widehat\phi(\xi)\widehat\mu(-\xi)\,\mathrm d\xi=\int\phi g\,\mathrm d x.
\]
The last equality is the bilinear form of Plancherel, obtained by conjugating one test. In particular the distribution defined by $g$ is positive. Its mollifications are nonnegative and converge to $g$ in $L^1$ locally, so $g$ is real and nonnegative almost everywhere. The uniqueness of Radon measures against smooth tests, using \cref{lem:mollifier-density,prop:continuous-density}, gives $\mu=g\,\mathrm d x$. Increasing compact cutoffs then give $\int g=\mu(\mathbb R^d)<\infty$. The norm equality is Plancherel.
\end{proof}

For $0<s<d$, the Riesz $s$-energy\index{energy!Riesz} of a finite positive measure is
\[
 I_s(\mu)=\iint|x-y|^{-s}\,\mathrm d\mu(x)\,\mathrm d\mu(y).
\]
This kernel has no fundamental-solution normalizing factor. We allow the value infinity.

\begin{theorem}\label{thm:riesz-fourier-energy}
For every finite positive measure on $\mathbb R^d$ and $0<s<d$,
\[
 I_s(\mu)=b_{d,s}\int_{\mathbb R^d}|\widehat\mu(\xi)|^2|\xi|^{s-d}\,\mathrm d\xi,
 \qquad b_{d,s}=\pi^{s-d/2}\frac{\Gamma((d-s)/2)}{\Gamma(s/2)}.
\]
Here $\Gamma(a)=\int_0^\infty v^{a-1}\mathrm e^{-v}\,\mathrm d v$ for $a>0$. The equality includes infinite energies.
\end{theorem}
\begin{proof}
The substitution $v=\pi t|z|^2$ gives
\[
 |z|^{-s}=\frac{\pi^{s/2}}{\Gamma(s/2)}\int_0^\infty t^{s/2-1}\mathrm e^{-\pi t|z|^2}\,\mathrm d t,
\]
including infinity at $z=0$. Gaussian inversion and Fubini against the finite product measure give
\[
 \iint\mathrm e^{-\pi t|x-y|^2}\,\mathrm d\mu(x)\,\mathrm d\mu(y)
 =t^{-d/2}\int\mathrm e^{-\pi|\xi|^2/t}|\widehat\mu(\xi)|^2\,\mathrm d\xi.
\]
The Fourier Gaussian is integrable, so that interchange is absolute. All terms after this identity are nonnegative. Tonelli therefore interchanges the $t$ and $\xi$ integrals even when the answer is infinite. For $\xi\neq0$, the substitution $v=\pi|\xi|^2/t$ evaluates the remaining integral as
\[
 \int_0^\infty t^{(s-d)/2-1}\mathrm e^{-\pi|\xi|^2/t}\,\mathrm d t
 =\pi^{(s-d)/2}|\xi|^{s-d}\Gamma((d-s)/2).
\]
Multiplication of the constants gives the formula. The single point $\xi=0$ has Lebesgue measure zero and does not affect it.
\end{proof}

\begin{example}\label{ex:gaussian-riesz-energy}
The probability measure $\mathrm d\mu=\mathrm e^{-\pi|x|^2}\,\mathrm d x$ on $\mathbb R^d$ has
\[
 I_s(\mu)=\left(\frac\pi2\right)^{s/2}\frac{\Gamma((d-s)/2)}{\Gamma(d/2)},\qquad 0<s<d.
\]
\end{example}
\begin{proof}
Polar integration of $\int\mathrm e^{-\pi|x|^2}\,\mathrm d x=1$ gives $s_{d-1}=2\pi^{d/2}/\Gamma(d/2)$. The transform in \cref{lem:fourier-gaussian} and the energy formula give
\[
 I_s(\mu)=b_{d,s}s_{d-1}\int_0^\infty\mathrm e^{-2\pi r^2}r^{s-1}\,\mathrm d r.
\]
The substitution $v=2\pi r^2$ evaluates the radial integral as $\Gamma(s/2)/(2(2\pi)^{s/2})$. Inserting the constants proves the formula.
\end{proof}

\section{Second derivatives and singular integrals}
\label{sec:hessian-singular-integrals}

Let $d\geq1$. Write $D^2v=(\partial_a\partial_bv)_{a,b=1}^d$, and use Euclidean norms for vectors and the Hilbert--Schmidt norm $|M|=(\sum_{a,b}|M_{ab}|^2)^{1/2}$ for matrices in $L^p$ norms. Constants $C_d$ and $C_{d,p}$ may change from line to line and depend only on the indicated parameters. A bound for a scalar Laplacian controls the full real Hessian in $L^p$, although it need not control that Hessian pointwise. Plancherel gives the estimate for $p=2$. The extension to other finite exponents uses cancellation away from the support of a function.

\begin{lemma}\label{lem:hessian-multiplier-kernel}
For $1\leq a,b\leq d$, define $T_{ab}$ on $L^2(\mathbb R^d)$ by
\[
 \widehat{T_{ab}f}(\xi)=-\frac{\xi_a\xi_b}{|\xi|^2}\widehat f(\xi),
 \qquad \xi\neq0,
\]
with multiplier zero at the origin. This is a self-adjoint operator of norm at most one. There is a smooth real kernel $K_{ab}$ away from zero such that
\[
 |K_{ab}(z)|\leq C_d|z|^{-d},\qquad
 |\nabla K_{ab}(z)|\leq C_d|z|^{-d-1}.
\]
If $f\in L^1\cap L^2$ has compact support, then almost everywhere outside that support
\[
 T_{ab}f(x)=\int K_{ab}(x-y)f(y)\,\mathrm d y.
\]
For $v\in C_c^\infty$, one has $\partial_a\partial_bv=T_{ab}(-\Delta v)$.
\end{lemma}
\begin{proof}
The multiplier is real and bounded by one in absolute value, so Plancherel gives the norm and adjoint assertions. The derivative identity follows by multiplying the transform of $-\Delta v$ by the displayed multiplier; both sides are $L^2$ functions.

\Needspace{16\baselineskip}
Put $H_t(z)=(4\pi t)^{-d/2}\mathrm e^{-|z|^2/(4t)}$. The Gaussian formula gives $\widehat H_t(\xi)=\mathrm e^{-4\pi^2t|\xi|^2}$. For $0<\varepsilon<R<\infty$, the operators
\[
 \int_\varepsilon^R\partial_a\partial_bH_t*f\,\mathrm d t
\]
have multiplier
\[
 -\frac{\xi_a\xi_b}{|\xi|^2}
    \bigl(\mathrm e^{-4\pi^2\varepsilon|\xi|^2}
          -\mathrm e^{-4\pi^2R|\xi|^2}\bigr)
\]
and converge to $T_{ab}f$ in $L^2$ as $\varepsilon\downarrow0$, $R\uparrow\infty$. Each finite integral is justified by the integrable smooth convolution kernels; bounded frequency multipliers and dominated convergence give the limit.

For $z\neq0$ define
\[
 K_{ab}(z)=\int_0^\infty\partial_a\partial_bH_t(z)\,\mathrm d t.
\]
Every differentiated integrand is bounded by a polynomial in $|z|/\sqrt t$ times $C t^{-(d+2+l)/2}\mathrm e^{-|z|^2/(4t)}$ for $l$ further derivatives. This is integrable uniformly on compact sets away from zero: the exponential controls $t\downarrow0$, and $t^{-(d+2+l)/2}$ controls $t\uparrow\infty$. Thus the kernel is smooth there. The change $t=r^2s$ shows $K_{ab}(rz)=r^{-d}K_{ab}(z)$ for $r>0$. Its bounded derivatives on the unit sphere give the two estimates.

On a compact set disjoint from the support of $f$, the integral of $|\partial_a\partial_bH_t(x-y)|$ in $t$ is at most $C_d|x-y|^{-d}$. Fubini and dominated convergence therefore identify the finite heat integrals with the asserted convolution in the limit, uniformly on that compact set. Their $L^2$ limit is $T_{ab}f$, proving the almost-everywhere identity. Exhaust the complement by such compact sets. This argument also applies for $d=1$, where $T_{11}=-1$ and the off-origin kernel vanishes.
\end{proof}

\begin{lemma}\label{lem:hessian-multiplier-weak-one}
For $f\in L^1(\mathbb R^d)\cap L^2(\mathbb R^d)$ and $\lambda>0$,
\[
 |\{x:|T_{ab}f(x)|>\lambda\}|\leq\frac{C_d}{\lambda}\|f\|_1.
\]
\end{lemma}
\begin{proof}
The dyadic cubes are $2^k(m+[0,1)^d)$, for $k\in\mathbb Z$ and $m\in\mathbb Z^d$. Choose the maximal cubes $Q$ for which the average of $|f|$ exceeds $\lambda$. Such a cube lies below a maximal one: the averages on its successive ancestors tend to zero, since their volumes tend to infinity and $f\in L^1$. Two maximal cubes are disjoint, and their parent averages give
\[
 \sum_Q|Q|\leq\lambda^{-1}\|f\|_1,\qquad
 \lambda<\frac1{|Q|}\int_Q|f|\leq2^d\lambda.
\]
Outside their union, $|f|\leq\lambda$ almost everywhere. Indeed a shrinking dyadic cube containing a Lebesgue point is contained in a centered cube of comparable volume; \cref{thm:lebesgue-differentiation} makes its average tend to $|f(x)|$.

Write $f_Q=|Q|^{-1}\int_Qf$, set $g=f$ off the selected cubes and $g=f_Q$ on each of them, and put $b_Q=(f-f_Q)\mathbf1_Q$. Then
\[
 \|g\|_1\leq\|f\|_1,\quad \|g\|_\infty\leq2^d\lambda,
 \quad \|g\|_2^2\leq2^d\lambda\|f\|_1,
 \qquad \sum_Q\|b_Q\|_1\leq2\|f\|_1.
\]
Each $b_Q$ has integral zero. Jensen's inequality and disjoint supports give $g\in L^2$ and convergence of $\sum_Qb_Q=f-g$ in $L^2$.

Let $c_Q$ be the center of $Q$ and let $Q^*$ be its concentric enlargement by the factor $4\sqrt d$. For $x\notin Q^*$, the segment from $x-c_Q$ to $x-y$, $y\in Q$, stays at distance at least $|x-c_Q|/2$ from zero. Cancellation and \cref{lem:hessian-multiplier-kernel} yield
\[
 \begin{aligned}
 |T_{ab}b_Q(x)|
 &=\left|\int_Q\bigl(K_{ab}(x-y)-K_{ab}(x-c_Q)\bigr)b_Q(y)\,\mathrm d y\right|\\
 &\leq C_d\frac{\operatorname{diam}Q}{|x-c_Q|^{d+1}}\|b_Q\|_1.
 \end{aligned}
\]
Polar integration outside a ball of radius comparable to $\operatorname{diam}Q$ gives
\[
 \int_{\mathbb R^d\setminus Q^*}|T_{ab}b_Q|\leq C_d\|b_Q\|_1.
\]
Put $B^*=\bigcup_QQ^*$. Its volume is at most $C_d\lambda^{-1}\|f\|_1$. Finite partial sums of the $b_Q$ converge in $L^2$, and so do their images under $T_{ab}$. Select an almost-everywhere convergent subsequence and use Fatou outside $B^*$ to get
\[
 \int_{\mathbb R^d\setminus B^*}|T_{ab}(f-g)|\leq C_d\|f\|_1.
\]
Chebyshev's inequality for this integral and the $L^2$ bound for $T_{ab}g$ bound the sets where the two summands exceed $\lambda/2$. Including $B^*$ proves the assertion.
\end{proof}
\index{Calder\'on--Zygmund decomposition}
\index{singular integral!weak $(1,1)$ estimate}

\begin{theorem}\label{thm:hessian-multiplier-lp}
For every $1<p<\infty$, $T_{ab}$ extends uniquely to a bounded operator on $L^p(\mathbb R^d)$. In particular
\begin{equation}\label{eq:euclidean-hessian-lp}
 \|D^2v\|_p\leq C_{d,p}\|\Delta v\|_p,
 \qquad v\in C_c^\infty(\mathbb R^d).
\end{equation}
\end{theorem}
\begin{proof}
Let $1<p<2$ and $f\in L^p\cap L^2$. For $\lambda>0$ split $f$ into $f\mathbf1_{\{|f|>\lambda\}}$ and $f\mathbf1_{\{|f|\leq\lambda\}}$. The first lies in $L^1\cap L^2$. Apply the weak estimate to it, and the $L^2$ estimate to the second, both with threshold $\lambda/2$. This gives
\[
 |\{|T_{ab}f|>\lambda\}|
 \leq\frac{C_d}{\lambda}\int_{|f|>\lambda}|f|
      +\frac{C_d}{\lambda^2}\int_{|f|\leq\lambda}|f|^2.
\]
For any measurable $h$, Tonelli applied to $|h(x)|^p=\int_0^{|h(x)|}p\lambda^{p-1}\,\mathrm d\lambda$ gives its norm as the integral of this distribution function, including infinite norms. Multiplying the inequality by $p\lambda^{p-1}$ and integrating therefore gives
\[
 \|T_{ab}f\|_p^p
 \leq C_dp\left(\frac1{p-1}+\frac1{2-p}\right)\|f\|_p^p.
\]
The two inner integrals are $\int_0^{|f|}\lambda^{p-2}\,\mathrm d\lambda$ and $\int_{|f|}^\infty\lambda^{p-3}\,\mathrm d\lambda$. Their stated values explain both restrictions on $p$. Plancherel handles $p=2$.

For $p>2$ put $p'=p/(p-1)<2$, take $f\in L^p\cap L^2$, and write $h=T_{ab}f$. For integers $N\geq1$, on $E_N=B(0,N)\cap\{|h|\leq N\}$ set $g_N=\mathbf1_{E_N}h|h|^{p-2}$. This belongs to $L^2\cap L^{p'}$. Self-adjointness, the proved $p'$ estimate and H\"older give
\[
 \int_{E_N}|h|^p=\int h\overline{g_N}
 =\int f\overline{T_{ab}g_N}
 \leq C_{d,p'}\|f\|_p\left(\int_{E_N}|h|^p\right)^{1/p'}.
\]
If the integral is nonzero, division bounds its $p$th root; if it is zero the bound holds directly. The sets $E_N$ increase to full measure, since $h\in L^2$. Monotone convergence proves the $p$ estimate. Density of $C_c^\infty$ in $L^p$, from \cref{lem:mollifier-density}, gives the unique extension. Finally apply the derivative identity in \cref{lem:hessian-multiplier-kernel} to each coefficient of $D^2v$.
\end{proof}
\index{singular integral!$L^p$ boundedness}
\index{Hessian!$L^p$ estimate}

For an integer $k\geq0$, $1\leq p<\infty$ and an open set $U\subset\mathbb R^d$, the Sobolev space\index{Sobolev space} $W^{k,p}(U)$ consists of $L^p$ functions whose distributional derivatives of order at most $k$ belong to $L^p$. Its norm is the sum of those derivative norms. On a compact smooth manifold, a finite coordinate cover and a smooth partition of unity define equivalent such norms. The chain rule and bounded derivatives of the coordinate changes prove equivalence for smooth functions; local smoothing and distributional passage to the limit give the same rule for weak derivatives.

\begin{lemma}\label{lem:gradient-lp-interpolation}
For $1\leq p<\infty$ and $\varepsilon>0$ there is $C_{d,\varepsilon}$ such that
\[
 \|\nabla v\|_p\leq\varepsilon\|D^2v\|_p+C_{d,\varepsilon}\|v\|_p,
 \qquad v\in C_c^\infty(\mathbb R^d).
\]
\end{lemma}
\begin{proof}
Taylor's formula along the $a$th coordinate direction gives, for $h>0$,
\[
 h\partial_av(x)=v(x+he_a)-v(x)
       -\int_0^h(h-t)\partial_a^2v(x+te_a)\,\mathrm d t.
\]
Translation invariance and the integral triangle inequality give $\|\partial_av\|_p\leq2h^{-1}\|v\|_p+(h/2)\|\partial_a^2v\|_p$. Sum over $a$ and choose $h$ sufficiently small in terms of $d,\varepsilon$.
\end{proof}

For $0<\alpha\leq1$, the norm $\|v\|_{C^{0,\alpha}(F)}$ on a compact set $F$ is $\sup_F|v|+\sup_{x\neq y\in F}|v(x)-v(y)|/|x-y|^\alpha$. The notation $C^{1,\alpha}$ requires continuously differentiable functions with first derivatives H\"older continuous with exponent $\alpha$.

\Needspace{10\baselineskip}
\begin{lemma}[Morrey's inequality]\label{lem:morrey-sobolev-representative}
Let $d<p<\infty$, $\alpha=1-d/p$, and $B'\Subset B$ be Euclidean balls. Every $v\in W^{1,p}(B)$ has a continuous representative with
\[
 \|v\|_{C^{0,\alpha}(\overline{B'})}\leq C\|v\|_{W^{1,p}(B)}.
\]
Here $C$ depends only on $d,p,B',B$. Consequently a $W^{2,p}_{\mathrm{loc}}$ function has a $C^{1,\alpha}_{\mathrm{loc}}$ representative.
\end{lemma}
\begin{proof}
First let $v$ be smooth. If $Q$ is a ball of radius $r$ and $x\in Q$, integrate $v(x)-v(z)$ along the segment from $x$ to $z$, then average in $z\in Q$. Polar coordinates about $x$ give
\[
 |v(x)-v_Q|\leq C_d\int_Q|\nabla v(z)|\,|x-z|^{1-d}\,\mathrm d z,
 \qquad v_Q=\frac1{|Q|}\int_Qv.
\]
To see the constant, write $z=x+s\theta$, whose allowed radii satisfy $0<s<R(\theta)\leq2r$ by convexity. Interchanging the segment integral and the radial integral produces the factor $\int_t^{R(\theta)}s^{d-1}\,\mathrm d s/|Q|\leq C_d$; converting back to volume gives $t^{1-d}$. H\"older now bounds the displayed integral by $C_{d,p}r^{1-d/p}\|\nabla v\|_{L^p(Q)}$. Indeed $(1-d)p'>-d$ exactly when $p>d$; the same calculation holds for $d=1$.

For $x,y\in B'$ with $r=|x-y|$ sufficiently small compared with $\operatorname{dist}(B',\partial B)$, use the common ball $Q=B(x,2r)\subset B$. Applying this bound at both points gives
\[
 |v(x)-v(y)|\leq C_{d,p}|x-y|^{1-d/p}\|\nabla v\|_{L^p(B)}.
\]
Using a fixed small ball about each $x\in B'$ also bounds $|v(x)|$ by $C\|v\|_{W^{1,p}(B)}$, since its average is bounded by $|Q|^{-1/p}\|v\|_p$. For more distant pairs this supremum bound gives the same H\"older estimate with a constant depending on $B',B$.

For a general $W^{1,p}$ function, mollify on a slightly smaller ball still containing $\overline{B'}$. The functions and their weak derivatives converge in $L^p$ by \cref{lem:mollifier-density,prop:distribution-smoothing}. The just proved estimate makes the approximations Cauchy uniformly and in the H\"older seminorm on $B'$. Their continuous limit agrees almost everywhere with $v$ and obeys the estimate. For $v\in W^{2,p}$ apply this argument to $v$ and each first derivative. Their smooth approximations converge locally uniformly; passing to the limit in the segment formula identifies the continuous derivative limits as the classical derivatives of the continuous limit. This gives the last assertion.
\end{proof}
\index{Morrey's inequality}
\index{Sobolev space!continuous representatives}


\section{Unitary invariant measures}
\label{sec:unitary-measures}
Identify $\mathbb C^n$ with $\mathbb R^{2n}$, and set $\langle z,w\rangle=\sum_jz_j\overline{w_j}$. Thus the real scalar product is $\operatorname{Re}\langle z,w\rangle$. The compact group $\mathrm U(n)$ consists of matrices $U$ with $U^*U=1$. Its invariant probability\index{measure!unitary invariant} has a direct Gaussian construction.

\Needspace{9\baselineskip}
\begin{proposition}\label{prop:unitary-invariant-probability}
There is a probability measure $\eta$ on $\mathrm U(n)$ invariant under left and right multiplication. Every compact metric space on which this group acts continuously and transitively has a unique invariant probability measure, obtained by applying a random matrix with law $\eta$ to one fixed point. In particular this defines invariant probabilities on the unit sphere, complex projective space, and every complex Grassmannian.
\end{proposition}
\begin{proof}
Let $X$ be a complex $n$ by $n$ matrix with density $\mathrm e^{-\pi\operatorname{tr}(X^*X)}$ for Lebesgue measure on its $2n^2$ real entries. The density integrates to one by the Gaussian formula. Singular matrices have measure zero. Indeed $|\det X|^2$ is a nonzero real polynomial; a nonzero polynomial has a null zero set by induction on the number of variables, viewing it as a polynomial in the last variable and applying Fubini off the null common zero set of its coefficients. Define
\[
 Q(X)=X(X^*X)^{-1/2}.
\]
The positive square root exists and is unique by the Hermitian spectral theorem. It depends continuously on positive definite matrices: along a convergent sequence, the roots are bounded, every convergent subsequence has square equal to the limiting matrix and is positive, and uniqueness identifies the limit; inversion is continuous away from zero eigenvalues. Thus $Q$ is Borel on the nonsingular matrices and is unitary.

For fixed $U,V\in\mathrm U(n)$, the matrix $UXV$ has the same Gaussian distribution. The map on entries is real orthogonal and preserves the displayed density. Functional calculus gives $Q(UXV)=UQ(X)V$. The distribution $\eta$ of $Q(X)$ is therefore invariant on both sides.

On a transitive compact space $Y$, let $\sigma$ be the distribution of $gy_0$ for $g$ with law $\eta$. Left invariance makes it invariant. For a continuous $f$ on $Y$, the number $\int f(gy)\,\mathrm d\eta(g)$ is independent of $y$, by transitivity and right invariance. Integrate this equality against any invariant probability on $Y$ and use Fubini. The result is $\int f\,\mathrm d\sigma$ for every such probability. To deduce uniqueness of the measures, approximate the indicator of any closed set $F\subset Y$ by the decreasing continuous functions $\max\{0,1-j\operatorname{dist}(y,F)\}$. Dominated convergence gives equality on closed sets; the $\pi$--$\lambda$ argument gives equality on their generated Borel $\sigma$-algebra.
\end{proof}

For an integer $0\leq k\leq n$, represent $\mathrm{Gr}_{\mathbb C}(k,n)$\index{Grassmannian!complex} by Hermitian matrices $P$ with $P^2=P$ and $\operatorname{tr}P=k$. These form a compact set, and $P$ is the orthogonal projection onto its image. The unitary action is $P\mapsto UPU^*$; choosing orthonormal bases proves transitivity. Projective space $\mathbb P^{n-1}(\mathbb C)$ is the case $k=1$. Its invariant probability will be denoted by $\sigma$. For unit representatives $e,e'$ of $\xi,\xi'$, set
\[
 d(\xi,\xi')=\arccos|\langle e,e'\rangle|\in[0,\pi/2].
\]
This is a metric: it is the minimum of the spherical distance $\delta(e,e')=\arccos\operatorname{Re}\langle e,e'\rangle$ between representatives on their two phase circles. To check the spherical triangle inequality, decompose two unit vectors relative to a third. Cauchy--Schwarz bounds their real scalar product below by $\cos(a+b)$ when the two intervening angles are $a,b$ and $a+b\leq\pi$; if $a+b>\pi$, the inequality is automatic. Choosing minimizing phases relative to one middle representative now gives the triangle inequality for $d$. It induces the quotient topology, by compactness and the continuous distance formula. The Hopf map\index{Hopf map} $q:S^{2n-1}\to\mathbb P^{n-1}(\mathbb C)$ sends a vector to its span.

\Needspace{9\baselineskip}
\begin{lemma}\label{lem:projective-ball-volume}
For $n\geq2$, $0<a\leq\pi/2$, and $\xi_0\in\mathbb P^{n-1}(\mathbb C)$,
\[
 \sigma\{\xi:d(\xi,\xi_0)<a\}=(\sin a)^{2n-2}.
\]
In particular these balls have measure at most $a^{2n-2}$.
\end{lemma}
\begin{proof}
Normalized surface measure on the sphere is unitary invariant, so its Hopf image is $\sigma$. The direction of a Gaussian vector $Z$ with density $\mathrm e^{-\pi|z|^2}$ has that sphere distribution, by polar coordinates. Take $\xi_0$ to be the first coordinate line. The independent variables $T_j=\pi|Z_j|^2$ have density $\mathrm e^{-t}$ on $t>0$, by planar polar coordinates. The event in question is
\[
 \frac{T_1}{T_1+\cdots+T_n}>\cos^2a,
 \quad\text{or}\quad T_1>\cot^2a\sum_{j=2}^nT_j.
\]
Conditioning on the last $n-1$ variables and integrating their densities gives $(1+\cot^2a)^{-(n-1)}=(\sin a)^{2n-2}$. Boundary equality has probability zero; the endpoint $a=\pi/2$ follows too. Unitary invariance gives every center.
\end{proof}

Write $\mathcal H=\mathrm{Gr}_{\mathbb C}(n-1,n)$, with invariant probability $\nu$. For $H\in\mathcal H$, let $m_H$ be its Euclidean Lebesgue measure in real dimension $2n-2$. Integrals depending on $H$ can be defined in Borel orthonormal coordinates. To construct these, project the standard coordinate vectors onto $H$ and apply Gram--Schmidt to a linearly independent choice. There are finitely many choices; their nonvanishing Gram determinants define open sets covering $\mathcal H$. Use the first successful choice to partition this cover into Borel sets. This gives a Borel isometry $T_H:\mathbb C^{n-1}\to H$, continuous on each underlying chart. Consequently $H\mapsto\int_H f(H,c)\,\mathrm d m_H(c)$ is Borel for nonnegative Borel $f$, by the parameter integral on fixed coordinates. A monotone-class argument from rectangles proves that parameter integral is Borel. The value is independent of the frame, since changes of orthonormal frame preserve Lebesgue measure.

\section{Circle selection and high-value sets}
\label{sec:high-value-stars}
The only selection needed below concerns closed subsets of circles with positive arc length. Their positive measure permits a direct construction that also proves measurability.

\Needspace{22\baselineskip}
\begin{lemma}\label{lem:positive-circle-selection}
\index{Borel selection!circle sections}Let $T$ be a separable metric space, and let $F\subset T\times[0,2\pi]$ be Borel, with every section $F_t$ closed and of positive Lebesgue measure. There is a Borel map $a:T\to[0,2\pi]$ with $a(t)\in F_t$.
\end{lemma}
\begin{proof}
Put $q(t)=|F_t|$ and $A_t(b)=|F_t\cap[0,b]|$. For fixed $b$, both functions are Borel in $t$, by a monotone-class argument starting with Borel rectangles. For each $t$, $A_t$ is continuous and nondecreasing, with values zero and $q(t)$ at the two endpoints. Define
\[
 a(t)=\inf\{b\in\mathbb Q\cap[0,2\pi]:A_t(b)\geq q(t)/2\}.
\]
Since $0<q(t)/2<q(t)$, the rational competitor set is nonempty, and continuity gives $A_t(a(t))=q(t)/2$. Countable infima give Borel measurability. If $a(t)\notin F_t$, an open interval about it misses $F_t$, so $A_t$ is constant there. A smaller rational point would be a competitor, contradicting the infimum. Hence $a(t)\in F_t$.
\end{proof}

\Needspace{14\baselineskip}
\begin{theorem}\label{thm:high-value-star-energy}
Let $n\geq2$, and let $u:D\to\mathbb R$ be Borel and subharmonic on every connected component of every affine complex-line section of the open set $D\subset\mathbb C^n$. Let $R>0$ satisfy $\overline{B(x,R)}\subset D$, and put $A=u(x)$. There is a Borel probability measure $\mu$ carried by
\[
 E=\{y:R/2\leq|y-x|\leq R,\ u(y)\geq A\}
\]
such that $I_{2n-2}(\mu)\leq C_nR^{-(2n-2)}$. The constant depends only on $n$.
\end{theorem}
\begin{proof}
For $r\in[R/2,R]$ and a projective direction $\xi$, choose a unit representative $e_0(\xi)$. It can be chosen Borel: in the rank-one projection matrix $P_\xi$, take the first index $j$ with $(P_\xi)_{jj}>0$ and put $e_0=P_\xi e_j/\sqrt{(P_\xi)_{jj}}$. This has norm one and a positive $j$th coordinate. The section
\[
 F_{r,\xi}=\{\theta\in[0,2\pi]:u(x+r\mathrm e^{\mathrm i\theta}e_0(\xi))\geq A\}
\]
is closed by upper semicontinuity on the complex line. It has positive length. Otherwise $u<A$ almost everywhere on the circle, while it is bounded above there, and its circle integral would be strictly less than $A$. This contradicts the submean inequality at $x$. The same inequality ensures that the negative part of that circle restriction is integrable, so the comparison of integrals is legitimate.

The total set of these sections is Borel, since $u$ and $e_0$ are Borel. Apply \cref{lem:positive-circle-selection} to choose a Borel phase, and let $e(r,\xi)$ be the resulting unit vector. Push normalized Lebesgue measure on $[R/2,R]$ times $\sigma$ forward by $\Psi(r,\xi)=x+re(r,\xi)$. This gives a Borel probability carried by $E$.

The energy estimate is independent of the phase choices. For unit representatives $e,e'$ of $\xi,\xi'$,
\[
 |re-r'e'|^2=(r-r')^2+rr'|e-e'|^2
 \geq(r-r')^2+\frac{R^2}{4}\bigl(2-2|\langle e,e'\rangle|\bigr).
\]
Writing $a=d(\xi,\xi')$ and using $\sin(a/2)\geq a/\pi$ gives a lower bound $(r-r')^2+R^2a^2/\pi^2$. Scale $r=Rt$, and give $[1/2,1]\times\mathbb P^{n-1}(\mathbb C)$ the metric $\rho^2=(t-t')^2+d(\xi,\xi')^2$ and product probability $m$. Then
\[
 |\Psi(Rt,\xi)-\Psi(Rt',\xi')|\geq cR\rho((t,\xi),(t',\xi'))
\]
for an absolute $c>0$. By \cref{lem:projective-ball-volume}, every $\rho$-ball of radius $a\leq1$ has $m$-measure at most $4a^{2n-1}$. A dyadic decomposition therefore gives, uniformly in its center $z$,
\[
 \int\rho(z,z')^{-(2n-2)}\,\mathrm d m(z')
 \leq1+\sum_{k=0}^\infty 2^{(k+1)(2n-2)}4\,2^{-k(2n-1)}<\infty.
\]
The ball bound also gives zero mass to a point, so the diagonal contributes no hidden atom. Integrating this bound in $z$ and using the distance inequality proves the asserted energy estimate.
\end{proof}

\Needspace{13\baselineskip}
\section{Projection densities and incidence}
\label{sec:complex-incidence}
\begin{lemma}\label{lem:complex-hyperplane-average}
For nonnegative Borel $f$ on $\mathbb C^n$,
\[
 \int_{\mathcal H}\int_H f(\xi)\,\mathrm d m_H(\xi)\,\mathrm d\nu(H)
 =a_n\int_{\mathbb C^n}f(\xi)|\xi|^{-2}\,\mathrm d\xi,
 \qquad a_n=\frac{s_{2n-3}}{s_{2n-1}}.
\]
\end{lemma}
\begin{proof}
Use polar coordinates in each $H$. Average its normalized sphere measure over $\nu$. The resulting probability on $S^{2n-1}$ is unitary invariant and hence is normalized full-sphere measure by \cref{prop:unitary-invariant-probability}. Thus the left side is
\[
 s_{2n-3}\int_0^\infty r^{2n-3}\int_{S^{2n-1}}f(r\theta)\,\frac{\mathrm d S(\theta)}{s_{2n-1}}\,\mathrm d r.
\]
Full-space polar coordinates give the right side. The frame construction just given and Tonelli justify all integrals for Borel $f$, including infinite values.
\end{proof}

\Needspace{13\baselineskip}
\begin{theorem}\label{thm:complex-projection-l2}
Let $\mu$ be a compactly supported finite positive measure on $\mathbb C^n$, $n\geq2$, with $I_{2n-2}(\mu)<\infty$. For almost every $H\in\mathcal H$, the orthogonal projection measure $(P_H)_*\mu$ has a density $g_H\in L^2(H,m_H)$, and
\[
 \int_{\mathcal H}\|g_H\|_2^2\,\mathrm d\nu(H)=\beta_n I_{2n-2}(\mu),
 \qquad \beta_n=\frac{a_n}{b_{2n,2n-2}}=\frac{\Gamma(n)}{\pi^{n-1}}>0.
\]
\end{theorem}
\begin{proof}
For $\xi\in H$, orthogonality gives $\operatorname{Re}\langle P_Hy,\xi\rangle=\operatorname{Re}\langle y,\xi\rangle$. Hence the Fourier transform of the projected measure is $\widehat\mu|_H$. Apply \cref{lem:complex-hyperplane-average} to $|\widehat\mu|^2$ and then \cref{thm:riesz-fourier-energy} to obtain
\[
 \int_{\mathcal H}\int_H|\widehat\mu(\xi)|^2\,\mathrm d m_H(\xi)\,\mathrm d\nu(H)
 =\beta_n I_{2n-2}(\mu)<\infty.
\]
The inner integral is a Borel function of $H$ in the fixed frames. It is finite almost everywhere. For those $H$, \cref{prop:measure-fourier-l2} gives the density and its norm. This proves the identity; the exceptional hyperplanes form a null set. The explicit constant uses the sphere area computed in \cref{ex:gaussian-riesz-energy} and $\Gamma(n)=(n-1)\Gamma(n-1)$, which follows by integration by parts. No jointly chosen density is required for the argument, since its norm equals the displayed Borel Fourier integral.
\end{proof}

\Needspace{9\baselineskip}
\begin{theorem}\label{thm:complex-line-incidence}
Let $K\subset\mathbb C^n$ be compact, with $n\geq2$, and let $E_j\subset K$, $j\geq1$, be Borel. Suppose probabilities $\mu_j$ satisfy $\mu_j(E_j)=1$ and $I_{2n-2}(\mu_j)\leq C<\infty$. Some affine complex line meets $E_j$ for infinitely many $j$.
\end{theorem}
\begin{proof}
Choose $M$ with $K\subset B(0,M)$. The compact space
\[
 \mathcal L_M=\{(H,c):H\in\mathcal H,\ c\in H,\ |c|\leq M\}
\]
parametrizes lines $c+H^\perp$ that can meet $K$. Give it the finite measure $\lambda$ defined by
\[
 \int h\,\mathrm d\lambda=\int_{\mathcal H}\int_{H\cap\overline{B(0,M)}}h(H,c)\,\mathrm d m_H(c)\,\mathrm d\nu(H)
\]
for nonnegative Borel $h$. In the Borel frames this is a product integral and defines a Borel measure; its mass is the volume of a $(2n-2)$-dimensional ball of radius $M$.

By inner regularity choose compact $F_j\subset E_j$ with $\mu_j(F_j)\geq1/2$. The incidence set
\[
 J_j=\{(H,c):c\in P_HF_j\}
\]
is compact: it is the image of $\mathcal H\times F_j$ under the continuous map $(H,y)\mapsto(H,P_Hy)$. For almost every $H$, \cref{thm:complex-projection-l2} and Cauchy--Schwarz give
\[
 \frac12\leq(P_H)_*\mu_j(P_HF_j)
 \leq\|g_{j,H}\|_2\,m_H(P_HF_j)^{1/2}.
\]
The first inequality follows since $F_j\subset P_H^{-1}(P_HF_j)$. Integrate in $H$ and apply Cauchy--Schwarz once more. Since $\nu$ is a probability,
\[
 \frac14\leq\lambda(J_j)\int_{\mathcal H}\|g_{j,H}\|_2^2\,\mathrm d\nu(H)
 \leq\lambda(J_j)\beta_n C.
\]
Thus $\lambda(J_j)\geq\delta=1/(4\beta_n C)>0$. The constant $C$ is positive because a probability of finite energy cannot have zero energy. The decreasing sets $\bigcup_{j\geq m}J_j$ have measure at least $\delta$. Since $\lambda$ is finite, continuity from above gives
\[
 \lambda\left(\bigcap_m\bigcup_{j\geq m}J_j\right)\geq\delta.
\]
Any pair in this set represents a line meeting infinitely many $F_j$, and hence infinitely many $E_j$.
\end{proof}

\Needspace{22\baselineskip}
\begin{corollary}\label{cor:borel-line-subharmonic-bound}
For $n\geq1$ and an open set $D\subset\mathbb C^n$, a finite Borel function $u:D\to\mathbb R$ that is subharmonic on every connected component of every affine complex-line section is locally bounded above.
\end{corollary}
\begin{proof}
For $n=1$ this is the upper semicontinuity of a planar subharmonic function. For $n\geq2$, otherwise there are $x_j\to p$ with $A_j=u(x_j)\to\infty$. Choose $R>0$ with $\overline{B(p,3R)}\subset D$ and discard finitely many terms so that $|x_j-p|<R/2$. By \cref{thm:high-value-star-energy}, the high-value annular sets about $x_j$ support probabilities of uniformly bounded $I_{2n-2}$-energy. They all lie in the fixed compact $\overline{B(p,2R)}$. By \cref{thm:complex-line-incidence}, one affine complex line meets infinitely many of them. Choose such points $y_j$ and a subsequence converging to $y$ on that line in the fixed compact set. A small disc in the line around $y$ lies in $D$, so eventually all the $y_j$ are in its single line-section component. Upper semicontinuity there contradicts $u(y_j)\geq A_j\to\infty$ and the finite value $u(y)$. This proves the local bound. Borel measurability is still an explicit hypothesis here; removing it belongs to the pointwise complex-line characterization.
\end{proof}
